\section{Булевы функции}
\label{sec:boolean-functions}

\subsection{Основные определения}

Булевой функцией от $n$ переменных называется отображение
\[
    f\colon \{0,1\}^{n}\to\{0,1\}.
\]
Множество всех булевых функций от $n$ переменных обозначают через $P_n$.
Функция может быть задана формулой, таблицей истинности, вектором значений или
одной из нормальных форм.

Переменная $x_i$ называется \emph{существенной} для функции $f$, если существуют два
набора, отличающиеся только значением $x_i$, на которых значения $f$ различны.
Если такого набора нет, переменная фиктивна.

К основным булевым операциям относятся
\[
\neg x,\qquad x\land y,\qquad x\lor y,\qquad
x\oplus y,\qquad x\to y=\neg x\lor y,
\]
а также штрих Шеффера $x\mid y=\neg(x\land y)$ и стрелка Пирса
$x\downarrow y=\neg(x\lor y)$.

Булева алгебра подчиняется законам коммутативности, ассоциативности,
дистрибутивности, поглощения и идемпотентности. Особенно важны законы де Моргана:
\[
    \neg(x\land y)=\neg x\lor\neg y,
    \qquad
    \neg(x\lor y)=\neg x\land\neg y.
\]

\subsection{Разложение по переменной и нормальные формы}

Для любой переменной $x_i$ справедливо разложение Шеннона:
\[
 f(x_1,\ldots,x_n)
 =\bigl(\neg x_i\land f|_{x_i=0}\bigr)
  \lor
  \bigl(x_i\land f|_{x_i=1}\bigr).
\]
Последовательное применение этой формулы даёт представление функции через
элементарные конъюнкции или дизъюнкции.

\emph{Минтермом} называется конъюнкция, содержащая каждую переменную ровно один раз
--- с отрицанием или без него. Минтерм равен единице ровно на одном наборе.
Дизъюнкция минтермов, соответствующих единичным строкам таблицы истинности,
называется совершенной дизъюнктивной нормальной формой (СДНФ):
\[
    f(x_1,\ldots,x_n)=
    \bigvee_{\alpha\colon f(\alpha)=1}
    x_1^{\alpha_1}\land\cdots\land x_n^{\alpha_n},
\]
где $x^1=x$, $x^0=\neg x$.

Двойственным образом строится совершенная конъюнктивная нормальная форма (СКНФ)
как конъюнкция макстермов, соответствующих нулевым строкам таблицы истинности.
Для функции, не равной тождественно нулю или единице, СДНФ и СКНФ определены
однозначно с точностью до перестановки членов.

\subsection{Полином Жегалкина}

Любая булева функция единственным образом представляется полиномом над полем
$\mathbb{F}_2$:
\[
 f(x_1,\ldots,x_n)=
 a_0\oplus
 \bigoplus_i a_i x_i\oplus
 \bigoplus_{i<j}a_{ij}x_ix_j\oplus\cdots\oplus
 a_{12\ldots n}x_1x_2\cdots x_n,
\]
где коэффициенты принимают значения $0$ или $1$. Это представление называется
полиномом Жегалкина, или алгебраической нормальной формой.

Функция называется \emph{линейной}, если её полином Жегалкина имеет степень не выше
одной:
\[
    f=a_0\oplus a_1x_1\oplus\cdots\oplus a_nx_n.
\]

\subsection{Двойственность, симметрия и монотонность}

Двойственной к $f$ называется функция
\[
    f^*(x_1,\ldots,x_n)=\neg f(\neg x_1,\ldots,\neg x_n).
\]
Замена $\land\leftrightarrow\lor$ и $0\leftrightarrow1$ в формуле приводит к
формуле для двойственной функции. Функция самодвойственна, если $f=f^*$.

Функция симметрична, если её значение не меняется при любой перестановке
переменных. Тогда значение зависит только от числа единиц в наборе.

На множестве $\{0,1\}^n$ вводится покоординатный порядок:
$\alpha\leq\beta$, если $\alpha_i\leq\beta_i$ для всех $i$.
Функция называется монотонной, если
\[
    \alpha\leq\beta\quad\Longrightarrow\quad f(\alpha)\leq f(\beta).
\]

\subsection{Функциональная полнота. Теорема Поста}

Система булевых функций $F$ называется \emph{полной}, если любую булеву функцию
можно выразить формулой, использующей функции из $F$ и переменные.
Классические полные системы:
\[
    \{\neg,\land\},\qquad \{\neg,\lor\},\qquad
    \{\mid\},\qquad \{\downarrow\}.
\]

В теореме Поста используются пять замкнутых классов:
\begin{enumerate}
    \item $T_0$ --- функции, сохраняющие ноль: $f(0,\ldots,0)=0$;
    \item $T_1$ --- функции, сохраняющие единицу: $f(1,\ldots,1)=1$;
    \item $S$ --- самодвойственные функции;
    \item $M$ --- монотонные функции;
    \item $L$ --- линейные функции.
\end{enumerate}

\textbf{Теорема Поста.}
Система $F$ функционально полна тогда и только тогда, когда она не содержится
целиком ни в одном из классов $T_0,T_1,S,M,L$. Иными словами, в системе должны
существовать функции, нарушающие каждое из пяти указанных свойств.

\subsection{Минимизация булевых формул}

Задача минимизации состоит в уменьшении числа литералов, конъюнкций или логических
операций при сохранении реализуемой функции. Основные методы:
\begin{itemize}
    \item карты Карно --- наглядный метод для малого числа переменных;
    \item метод Квайна--Мак-Класки --- табличное построение простых импликант;
    \item использование законов булевой алгебры и поглощения.
\end{itemize}

Сокращённая ДНФ строится из простых импликант; минимальная ДНФ выбирается среди
покрытий всех единичных наборов минимальной стоимости.

\subsection{Что важно сказать на экзамене}

Булева функция задаётся таблицей истинности и имеет канонические представления:
СДНФ, СКНФ и полином Жегалкина. Полнота системы функций проверяется по теореме
Поста. Для практической реализации логических схем существенны минимизация формул
и выбор базиса.

\newpage